\chapter{獣害対策における監視の課題}
\label{chap:research-background}

獣害対策の監視では，動物の出現状況を把握し，必要な情報を継続的に得ることが課題となる．
本章では，監視情報の利用目的と，現地での画像収集，通信および電源に関する制約を整理する．
装置や検出モデルの基礎は第\ref{chap:background}章で述べる．

\section{動物の出現情報と対策の関係}
動物の出現時刻と場所を記録することは，見回りや対策箇所の検討を支援する．
トレイルカメラは人が常駐しない地点でも画像を取得できるため，出現状況を継続して調べる手段となる\cite{Beery2019,IEICE2020}．
ただし，画像による監視は情報を得る手段であり，その導入だけで農作物被害の減少を保証するものではない．
\writehere{研究対象地域の被害状況，対象動物，情報を利用する主体と対応方法を，確認した資料に基づいて補う．}

\section{画像回収と確認作業の負担}
画像を現地で回収する運用では，出現情報を確認するまでに回収待ちの時間が生じる．
監視地点の増加は移動と画像確認の負担を増やし，動物が写っていない画像の選別も必要となる\cite{Beery2019}．
このため，画像を遠隔へ転送する仕組みと，確認すべき画像を選ぶ仕組みを組み合わせる必要がある．

\section{屋外設置における通信と電源の制約}
複数地点での監視では，通信経路と電源を各地点に確保することが運用上の制約となる．
画像認識と通知を担う機器を電池で動かす場合には，処理中の電力だけでなく待機中の電力も考慮する必要がある\cite{IEICE2020,Sato2022}．
通信回線や計算機を共有する構成では，各地点に配置する機器を簡素化できる一方，共有地点までの通信到達性と共有機器の継続稼働が必要となる．

\section{本研究で取り組む課題}
本研究は，複数の撮影地点から画像を集約し，動物画像の見逃しを考慮しながら広域通信へ送る画像を選別する課題に取り組む．
撮影を担うカメラ，一次選別を担う共有エッジサーバ，再判定と通知を担うクラウドに処理を分ける．
この構成について，画像の選別性能，通信，通知時間および電源条件を対応付けて検討する．
